\documentclass[11pt]{article}
\usepackage[a4paper,margin=27mm]{geometry}
\usepackage{amsmath,amssymb,amsthm,mathtools}
\usepackage{bm}
\usepackage{microtype}
\usepackage{booktabs}
\usepackage[numbers,sort&compress]{natbib}
\usepackage[colorlinks=true,allcolors=blue]{hyperref}
\hypersetup{
  pdftitle={A two-use history advantage under one-bit coarse-grained environment feedback},
  pdfauthor={Xiaojian Zeng}
}

\newtheorem{theorem}{Theorem}
\newtheorem{lemma}{Lemma}
\newcommand{\cI}{\mathcal I}
\newcommand{\cR}{\mathcal R}
\newcommand{\cT}{\mathcal T}
\newcommand{\id}{\mathrm{id}}
\newcommand{\Tr}{\operatorname{Tr}}
\newcommand{\diag}{\operatorname{diag}}

\title{A two-use history advantage under one-bit coarse-grained environment feedback}
\author{Xiaojian Zeng\\
\small PBC School of Finance, Tsinghua University, Beijing 100190, China\\
\small \texttt{xiaojianzeng@foxmail.com}}
\date{}

\begin{document}
\maketitle

\begin{abstract}
Environment-assisted recovery is usually formulated with every measured Kraus outcome
resolved.  We instead disclose one bit obtained by pairing four outcomes of a full-rank
qubit channel.  For two uses of an explicit generalized-amplitude-damping instrument,
we prove that every controller whose recovery sees only the current bit has
entanglement fidelity at most $456/625$.  A controller retaining the preceding bit
attains $1831/2500$, giving a certified advantage of at least $7/2500$.  Since the two
policy classes coincide for one use, two uses are the earliest possible depth at which
coarse-record history can help.  The upper bound permits arbitrary qubit CPTP
recoveries.
\end{abstract}

\noindent\textbf{Keywords:} environment-assisted correction; quantum feedback;
coarse-grained measurement; entanglement fidelity; generalized amplitude damping

\section{Introduction}

Classical information obtained from an environment can guide recovery of a noisy
quantum system.  In the resolved-record setting, each measured outcome labels a single
Kraus operator.  Gregoratti and Werner derived optimal discrete-time recovery formulas
for this setting~\cite{GregorattiWerner2004}.  Rexiti, Oskouei, and Mancini subsequently
compared current-outcome and history-conditioned discrete feedback and showed that the
two have equal performance for their rank-two qubit protocol~\cite{RexitiEtAl2024}.
Related Bayesian--Markovian distinctions also arise in continuous monitored feedback,
but with a different dynamics, observation process, and objective~\cite{WisemanEtAl2002}.

Finite-bandwidth readout changes the mathematical object presented to the controller.
If several orthogonal environment outcomes are assigned the same classical symbol, an
observed branch is a multi-Kraus completely positive operation; the unobserved
alternatives must be summed incoherently before control.  A polar correction for each
hidden Kraus operator is then unavailable.  We ask the first depth at which retaining
the history of such a coarse record can improve entanglement fidelity over all
current-bit-only recovery policies.

Our contribution is an exact two-use separation.  We give a full-Choi-rank generalized
amplitude-damping (GAD) channel and a fixed balanced pairing of its four environment
outcomes.  A trace-norm data-processing argument bounds each Bloch column attainable
after one current-bit correction.  Exact factorisation of a two-use current-bit policy,
followed by Frobenius Cauchy--Schwarz, yields a global bound over arbitrary qubit CPTP
recoveries.  An explicit history policy exceeds it.  The conclusion concerns the
information structure of this fixed finite-dimensional instrument; it is not an
experimental robustness or continuous-time statement.

\section{Model and policy classes}

Let the system be $\mathbb C^2$.  A fixed minimal Kraus representation
$\bm K=(K_0,K_1,K_2,K_3)$ defines
\begin{equation}
  \cT(\rho)=\sum_{x=0}^{3}K_x\rho K_x^\dagger,
  \qquad \sum_xK_x^\dagger K_x=I.
  \label{eq:channel}
\end{equation}
The representation is physically tied to an orthogonal measurement of a fresh
four-dimensional environment.  A balanced map $g:\{0,1,2,3\}\to\{0,1\}$ discloses
only $b=g(x)$, so the observed instrument is
\begin{equation}
  \cI_b(\rho)=\sum_{x:g(x)=b}K_x\rho K_x^\dagger.
  \label{eq:instrument}
\end{equation}
The fine outcome $x$ is never available.  Independent fresh environments and the same
instrument are used at each step.

After record $b_{1:k}$, a deterministic policy applies a qubit CPTP map
$\cR_{k,b_{1:k}}$.  Its unnormalised branches obey
\begin{equation}
 \Gamma_{\varnothing}=\id,\qquad
 \Gamma_{b_{1:k}}=\cR_{k,b_{1:k}}\circ\cI_{b_k}
                     \circ\Gamma_{b_{1:k-1}},
 \qquad
 \Phi_n=\sum_{b_{1:n}}\Gamma_{b_{1:n}}.
 \label{eq:branches}
\end{equation}
The current-bit class allows time-dependent $\cR_{k,b_k}$, whereas the full-history
class allows $\cR_{k,b_{1:k}}$.  Both use the current bit immediately.  Neither class
has access to hidden fine outcomes, postselection, or persistent quantum memory.
Randomised recovery is already included through CPTP mixing.

Performance is the entanglement fidelity for the maximally mixed qubit,
\begin{equation}
 F_e(\Phi)=\langle\Omega| (\id\otimes\Phi)
 (|\Omega\rangle\langle\Omega|)|\Omega\rangle,
 \qquad |\Omega\rangle=(|00\rangle+|11\rangle)/\sqrt2.
 \label{eq:fidelity}
\end{equation}
Let $F_n^M$ and $F_n^B$ be the maxima in the current-bit and full-history classes and
$\Delta_n=F_n^B-F_n^M$.  Compactness of the finite products of qubit-CPTP Choi sets
ensures that these maxima exist.  At $n=1$ the two policy classes are identical.

\section{A full-rank coarse GAD instrument}

Set
\begin{equation}
 p=\frac{18}{25},\qquad q=1-p=\frac7{25},\qquad
 \gamma=\frac3{10},\qquad s=\sqrt{1-\gamma}=\sqrt{\frac7{10}},
 \label{eq:parameters}
\end{equation}
and use the ordered GAD Kraus operators
\begin{equation}
 \begin{aligned}
 K_0&=\sqrt p\,\diag(1,s),\qquad&
 K_1&=\sqrt{p\gamma}\,|0\rangle\langle1|,\\
 K_2&=\sqrt q\,\diag(s,1),\qquad&
 K_3&=\sqrt{q\gamma}\,|1\rangle\langle0|.
 \end{aligned}
 \label{eq:kraus}
\end{equation}
Equation~\eqref{eq:channel} follows by direct diagonal summation.  The off-diagonal
operators $K_1$ and $K_3$ occupy distinct nonzero matrix units.  Within the diagonal
subspace, the determinant of the two coordinate vectors from $K_0$ and $K_2$ is
$\sqrt{pq}(1-s^2)=\sqrt{pq}\,\gamma$, whose squared modulus is
\begin{equation}
 pq\gamma^2=\frac{567}{31250}>0.
\end{equation}
Thus the four operators are linearly independent and the channel has Choi rank four.

We disclose the pairing $03|12$:
\begin{equation}
 \cI_0(\rho)=K_0\rho K_0^\dagger+K_3\rho K_3^\dagger,
 \qquad
 \cI_1(\rho)=K_1\rho K_1^\dagger+K_2\rho K_2^\dagger.
 \label{eq:pairing}
\end{equation}
Both observed operations genuinely contain two hidden alternatives.

\section{Exact separation}

\begin{theorem}\label{thm:main}
For the instrument in Eqs.~\eqref{eq:parameters}--\eqref{eq:pairing}, every two-use
current-bit policy with arbitrary qubit CPTP recoveries satisfies
\begin{equation}
 F_2^M\le\frac{456}{625}.
 \label{eq:markovbound}
\end{equation}
There is a full-history policy with fidelity $1831/2500$.  Consequently
$\Delta_2\ge 7/2500>0$, and the earliest use count at which a coarse-record history
advantage can occur is $n_*=2$.
\end{theorem}

\subsection{A feasible full-history policy}

Let $X$ be the Pauli bit flip and choose
\begin{equation}
 R_{1,0}=I,\quad R_{1,1}=X,\qquad
 R_{2,b_1b_2}=\begin{cases}I,&b_1=0,\\X,&b_1=1.\end{cases}
 \label{eq:historypolicy}
\end{equation}
These recoveries are unitary and hence CPTP.  Keeping the fine outcomes only as an
algebraic device for evaluating the incoherent sums, the entanglement fidelity is
\begin{equation}
 F_e=\frac14\sum_{x_1,x_2=0}^{3}
 \left|\Tr\!\left[R_{2,g(x_1)g(x_2)}K_{x_2}
 R_{1,g(x_1)}K_{x_1}\right]\right|^2.
 \label{eq:pathformula}
\end{equation}
Direct multiplication of the sixteen paths leaves the following multiset of nonzero
trace amplitudes:
\begin{equation}
 p(2-\gamma),\quad 2\sqrt{pq(1-\gamma)},\quad
 \gamma\sqrt{pq},\quad p\gamma,\quad
 \sqrt{pq}(2-\gamma),\quad 2q\sqrt{1-\gamma}.
 \label{eq:amplitudes}
\end{equation}
Substitution into Eq.~\eqref{eq:pathformula} gives the exact lower bound
\begin{equation}
 F_2^B\ge F_{\mathrm{hist}}
 =1-\gamma+\frac{p\gamma^2}{2}
 =\frac{1831}{2500}.
 \label{eq:historyvalue}
\end{equation}

\subsection{A global current-bit upper bound}

Consider one stage with arbitrary qubit CPTP maps $R_0,R_1$ and define
\begin{equation}
 C=R_0\circ\cI_0+R_1\circ\cI_1.
 \label{eq:onestage}
\end{equation}
This averaged map is CPTP.  Write its affine Bloch action as
$\bm r\mapsto\bm t+T\bm r$, equivalently
$C(\sigma_j)=\sum_iT_{ij}\sigma_i$ for Pauli matrices $\sigma_j$.

Every positive trace-preserving map contracts the trace norm on Hermitian inputs.
Indeed, the adjoint is positive and unital and therefore preserves the operator order
interval $[-I,I]$ used in the variational characterisation of the Hermitian trace norm.
Triangle inequality in Eq.~\eqref{eq:onestage} then yields
\begin{equation}
 2\|T\bm e_j\|_2=\|C(\sigma_j)\|_1
 \le \sum_{b=0}^{1}\|\cI_b(\sigma_j)\|_1.
 \label{eq:contraction}
\end{equation}
The equality on the left uses trace preservation: $C(\sigma_j)$ is a traceless
Hermitian qubit operator.

For the instrument~\eqref{eq:pairing}, direct evaluation gives
\begin{align}
 \cI_0(X)&=psX,& \cI_1(X)&=qsX,\nonumber\\
 \cI_0(Y)&=psY,& \cI_1(Y)&=qsY,\label{eq:pauliimages}\\
 \cI_0(Z)&=\diag(p,\gamma-p),&
 \cI_1(Z)&=\diag(q-\gamma,-q).\nonumber
\end{align}
Because $q\le\gamma\le p$, the two $Z$-image trace norms sum to $2p$.
Equation~\eqref{eq:contraction} therefore bounds the three Bloch columns by
\begin{equation}
 \|T\bm e_x\|_2\le s,\qquad
 \|T\bm e_y\|_2\le s,\qquad
 \|T\bm e_z\|_2\le p,
\end{equation}
and hence
\begin{equation}
 \|T\|_F^2\le 2(1-\gamma)+p^2=\frac{1199}{625}.
 \label{eq:frobeniuscap}
\end{equation}

For a two-use current-bit policy, summing both records factorises exactly:
\begin{equation}
 \Phi_2^M=C_2\circ C_1,\qquad
 C_k=\sum_{b=0}^{1}R_{k,b}\circ\cI_b.
 \label{eq:factorisation}
\end{equation}
No unitality or unitary restriction is used here.  For trace-preserving qubit maps with
Bloch blocks $T_1,T_2$, the normalised entanglement fidelity of the composition is
\begin{equation}
 F_e(C_2\circ C_1)=\frac{1+\Tr(T_2T_1)}4.
 \label{eq:blochfidelity}
\end{equation}
The affine translation does not contribute to the superoperator trace underlying
Eq.~\eqref{eq:blochfidelity}.  Frobenius Cauchy--Schwarz and
Eq.~\eqref{eq:frobeniuscap} now give
\begin{equation}
 \Tr(T_2T_1)\le\|T_2\|_F\|T_1\|_F\le\frac{1199}{625},
\end{equation}
which proves Eq.~\eqref{eq:markovbound}.  Combining the global upper bound with
Eq.~\eqref{eq:historyvalue} yields
\begin{equation}
 \Delta_2\ge\frac{1831}{2500}-\frac{456}{625}
 =\frac7{2500}>0.
\end{equation}
At one use, a record history contains only the current bit, so the two policy classes
coincide and $\Delta_1=0$.  This completes the proof of Theorem~\ref{thm:main}.

\section{Interpretation and limitations}

The theorem shows that an environment bit available for immediate recovery need not
remain adequate for the next recovery: retaining the earlier coarse bit can improve
the two-use optimum.  The mechanism is not Hilbert-space size alone.  Coarse graining
turns each disclosed outcome into a two-Kraus operation, removing the resolved record
on which the existing qubit no-advantage argument relies~\cite{GregorattiWerner2004,RexitiEtAl2024}.
The separation remains valid even when every current-bit recovery is an arbitrary
qubit CPTP map.

Several conclusions are deliberately excluded.  The proof does not determine either
exact optimum $F_2^M$ or $F_2^B$; its $7/2500$ is a certified lower bound on their
difference, not the exact gap.  It does not optimise over instruments or establish a
smallest positive gap on the continuous rank-four domain.  It supplies no robustness
radius for perturbations or laboratory errors and makes no claim about an $n>2$
optimum.  Finally, the literature comparison is bounded rather than an exhaustive
priority search; the relation to broader partially observed quantum control should be
reviewed by the author before submission.

\section*{Reproducibility}

The exact witness, Pauli images, rank certificate, trace-norm bounds, hidden-path
calculation, and an independent Liouville evaluation are implemented in the accompanying
Python script.  The calculations use rational and symbolic algebra for the logical
certificate; floating-point output is only a replay.  The accompanying reproducibility
file records the environment and invocation.

\section*{Statements and declarations}

\paragraph{Funding.} No funding information was supplied for this draft; the author
must confirm or replace this statement before submission.

\paragraph{Competing interests.} No competing-interest information was supplied for
this draft; the author must provide the final declaration before submission.

\paragraph{Data and code availability.} No empirical data were used.  The exact input
record and verification code are included with the local manuscript package.

\paragraph{AI use declaration.} OpenAI \texttt{gpt-5.6-sol} was used on
13--14 September 2026 to assist problem formulation, bounded literature discovery,
derivation, code generation, exact and numerical cross-checks, and manuscript drafting.
All outputs remain subject to human verification; the exact calculations were also
replayed in independent representations against the stated equations and primary-source
records.  AI is not an author and cannot assume accountability.  This local draft
requires independent review, correction, and approval by the named human author before
any submission.

\bibliographystyle{unsrtnat}
\bibliography{references}

\end{document}
